\documentclass{article}
\usepackage[UTF8]{ctex}
\usepackage{amsmath,amssymb,amsfonts}
\usepackage{enumitem}
\usepackage{geometry}
\geometry{a4paper, margin=1in}
% === 错题本排版优化（不参与题目标识） ===
\setlength{\parindent}{0pt}
\setlength{\parskip}{0.35em plus 0.10em minus 0.08em}
\setlength{\emergencystretch}{3em}
\setlist{itemsep=0.25em, topsep=0.40em, parsep=0pt}
\allowdisplaybreaks[2]
% === 错题本排版优化结束 ===
\begin{document}
\section*{武汉理工大学2025年高等代数真题}
\section*{127．武汉理工大学 2025 年高等代数试卷}

\subsection{题目1}
1．已知矩阵 \(A=\left(\begin{array}{cccc}2 a_{1} b_{1} & a_{1} b_{2}+a_{2} b_{1} & \cdots & a_{1} b_{n}+a_{n} b_{1} \\ a_{2} b_{1}+a_{1} b_{2} & 2 a_{2} b_{2} & \cdots & a_{2} b_{n}+a_{n} b_{2} \\ \vdots & \vdots & & \vdots \\ a_{n} b_{1}+a_{1} b_{n} & a_{n} b_{2}+a_{2} b_{n} & \cdots & 2 a_{n} b_{n}\end{array}\right)\) ，求 \(|A|\) ．

\subsection{题目2}
2．已知矩阵

\[
A^{\bullet}=\left(\begin{array}{lll}
5 & 1 & 0 \\
1 & 1 & 0 \\
0 & 0 & 1
\end{array}\right) .
\]

存在矩阵 \(B\) 满足 \(A^{*} B A=44 E-B A\) ，其中 \(A^{*}\) 为 \(A\) 的伴随矩阵，\(E\) 为单位矩阵。求 \(B\) 。

\subsection{题目3}
3．已知向量组 \(\alpha_{1}=(1+a, 1,1,1), \alpha_{2}=(2,2+a, 2,2), \alpha_{3}=(3,3,3+a, 3), \alpha_{4}=(4,4,4,4+a)\) 。\\
（1）\(a\) 为何值时，向量组 \(\alpha_{1}, \alpha_{2}, \alpha_{3}, \alpha_{4}\) 线性相关？\\
（2）在向量组 \(\alpha_{1}, \alpha_{2}, \alpha_{3}, \alpha_{4}\) 线性相关的前提下，求 \(\alpha_{1}, \alpha_{2}, \alpha_{3}, \alpha_{4}\) 的一个极大线性无关组，并用其线性表出其余向量．

\subsection{题目4}
4．已知 \(A\) 为 3 阶方阵，满足 \(|A-E|=|A-2 E|=|A+E|=\lambda\) 。

\begin{itemize}
  \item[（1）] 当 \(\lambda=0\) 时，求行列式 \(|A+3 E|\) 。
  \item[（2）] 当 \(\lambda=2\) 时，求行列式 \(|A+3 E|\) 。
\end{itemize}

\subsection{题目5}
5．设集合 \(W=\left\{f(t)=a_{0}+a_{1} t+a_{2} t^{2}+a_{3} t^{3} \mid f(2)=0, a_{0}, a_{1}, a_{2}, a_{3} \in \mathbb{R}\right\}\) ．

\begin{itemize}
  \item[（1）] 证明：\(W\) 是一个线性空间．
  \item[（2）] 求 \(W\) 的一组基和维数。
\end{itemize}

\subsection{题目6}
6．设 \(\varepsilon_{1}, \varepsilon_{2}, \varepsilon_{3}\) 为欧氏空间 \(V\) 的一组标准正交基，设向量组 \(\alpha_{1}, \alpha_{2}, \alpha_{3}\) 为 \(V\) 的另一组基，其中

\[
\alpha_{1}=3 \varepsilon_{1}+2 \varepsilon_{2}+\varepsilon_{3}, \alpha_{2}=2 \varepsilon_{1}+\varepsilon_{2}, \alpha_{3}=\varepsilon_{1} .
\]

设 \(T\) 为 \(V\) 上的一个线性变换，且满足 \(T \varepsilon_{1}=\alpha_{1}, T \varepsilon_{2}=\alpha_{1}+\alpha_{2}, T \varepsilon_{3}=\alpha_{1}+\alpha_{2}+\alpha_{3}\) ．

\begin{itemize}
  \item[（1）] 判断 \(T\) 是否为可逆变换？是否为正交变换？
  \item[（2）] 求 \(2 T-T^{-1}\) 在 \(\alpha_{1}, \alpha_{2}, \alpha_{3}\) 这组基下的矩阵。
\end{itemize}

\subsection{题目7}
7．设非齐次线性方程组 \(A X=b\) ：

\[
\left\{\begin{array}{l}
x_{3}=1 \\
x_{1}+x_{2}+x_{3}=1 \\
x_{1}+x_{2}=1
\end{array}\right.
\]

定义欧氏空间 \(\mathbb{R}^{3}\) 中的内积：设 \(X, Y \in \mathbb{R}^{3}, X=\left(x_{1}, x_{2}, x_{3}\right)^{\prime}, Y=\left(y_{1}, y_{2}, y_{3}\right)^{\prime}\) ，

\[
(X, Y)=2 x_{1} y_{1}+x_{2} y_{2}+x_{3} y_{3} .
\]

范数为 \(\|X\|=\sqrt{(X, X)}\) ．

\begin{itemize}
  \item[（1）] 求集合 \(W=\left\{X \in \mathbb{R}^{3} \mid\|A X-b\|=\min _{Y \in \mathbb{R}^{3}}\|A Y-b\|\right.\) 。
  \item[（2）] 求满足 \(\left\|X_{0}\right\|=\max _{X \in W}\|X\|\) 的 \(X_{0}\) ．
\end{itemize}

\subsection{题目8}
8．设 \(A=\alpha \beta^{\prime}\) ，其中 \(\alpha, \beta\) 为复数域上的 \(n\) 维列向量，记 \(k=\alpha^{\prime} \beta\) 。

\begin{itemize}
  \item[（1）] 求 \(A\) 的最小多项式。
  \item[（2）] 求 \(A\) 的 Jordan 标准形．
\end{itemize}

\subsection{题目9}
9．设 \(A\) 为 \(n\) 阶方阵，若 \(A+E\) 可逆，设 \(f(A)=(E-A)(E+A)^{-1}\) 。

\begin{itemize}
  \item[（1）] 证明：\((E+f(A))(E+A)=2 E\) 。
  \item[（2）] 证明：\(f(f(A))=A\) ．
\end{itemize}
\end{document}
